\section{Backgrounds}
\label{sec:intro_background}


Modern neural networks are very powerful predictive models, but they are often incapable of recognizing when their predictions may be wrong.


Uncertainty Quantification (UQ) is a key building block of reliable machine learning systems.

\begin{table}[h]
\centering
\begin{tabular}{ll}
\hline
\textbf{Notation} & \textbf{Meaning} \\
\hline
$P, p$ & probability measure, density or mass function \\
$\mathcal{X}, x, x_i$ & instance space, instance \\
$\mathcal{Y}, y, y_i$ & output space, outcome \\
$\gF, f, f_\theta$ & function space, prediction model, model parametrized by $\theta$ \\
%$\pi(y)$ & possibility/plausibility of outcome $y$ \\
$f(x)$ & outcome $y \in \mathcal{Y}$ for instance $x$ predicted by model $f$ \\
$p_{\theta}(y \mid x) = p(y \mid x, \theta)$ & probability of  $y$ given $x$, prediction model parametrized by  $\theta$ \\
$\ell$ & loss function \\
%$h^*, \hat{h}$ & risk-minimizing hypothesis, empirical risk minimizer \\
$\mathbb{E}(X)$ & expected value of random variable $X$ \\
$\text{Var}(X)$ & variance of random variable $X$ \\
$\mathbbm{1} \{\cdot\}$ & indicator function \\
$\gD, \Dtrain, \Dcal$ & dataset, training dataset, calibration dataset \\ 
\hline
\end{tabular}
%\vspace{1em}
\caption{Notations used through this paper.}
\end{table}

\subsection{Aleatoric and Epistemic Uncertainty}

A useful framework for analyzing uncertainty in machine learning prediction models is the distinction of \textit{aleatoric} and \textit{epistemic} uncertainty.

\textit{Aleatoric uncertainty} (from Latin ``alea" meaning dice) represents irreducible environmental stochasticity—inherent randomness in states and actions that cannot be reduced through additional data collection.

\textit{Epistemic uncertainty} (from Greek ``episteme" meaning knowledge) captures model uncertainty due to limited information during learning, which can be reduced with more data and better model understanding.

The standard decomposition of predictive uncertainty into aleatoric and epistemic parts is usually expressed via the law of total variance as in \Eqref{eq:aleatoric_epistemic_decomp} \cite{mackay2003information}, where $x\in \gX$ is the input, $y\in \gY$ is the output, $\gD$ is the dataset and $\theta$ is the learned model parameters.

\begin{equation}
\text{Var}[y|x, \gD] = \underbrace{\mathbb{E}_{p(\theta|D)}[\text{Var}[y|x,\theta]]}_{\text{Aleatoric}} + \underbrace{\text{Var}_{p(\theta|\gD)}[\mathbb{E}[y|x,\theta]]}_{\text{Epistemic}}
\label{eq:aleatoric_epistemic_decomp}
\end{equation}



In practice, disentangling the two types of uncertainty is nontrivial. \cite{Gal2017uncertainty}  models aleatoric uncertainty through learned variance $\delta^2(x)$ as network output, while epistemic uncertainty emerges through Monte Carlo Dropout variance across multiple forward passes.
\cite{devries2018learning} measures the total uncertainty and the aleatoric uncertainty and consider the epistemic uncertainty as the difference between the two.

\subsection{Conformal Prediction}
 Conformal prediction (CP) is a popular distribution-free UQ method that provides finite sample coverage guarantees without placing assumptions on the underlying data generating process \cite{vovk2005algorithmic}.
